% n3_trial_01_trap_depth — Unified trap-depth picture (N=3 LLM PNG → TikZ trial #1)
% Authored from briefing.md only; topic-related fixtures are intentionally not consulted.
% Layout: left column = 3 numbered rows (Experiment / Math / Molecular), center = green
% "Convergence" connector, right = BandDiagram + g(E_t) side lobes.

\documentclass[border=2pt]{standalone}
\usepackage{polymer-paper-preamble}

% Convergence-connector accent (briefing §5: green outlined arrow/box).
% cTeal is the closest palette neighbor to "green"; we lean green-ward via mix with black.
\colorlet{cConv}{cTeal!75!black}

\begin{document}
\begin{tikzpicture}[x=1cm, y=1cm,
  rowFrame/.style={draw=cTeal!55, line width=0.6pt, fill=cTeal!7,
                   rounded corners=2.5mm},
  numCircle/.style={circle, draw=cConv, fill=cConv, text=white,
                    inner sep=0pt, minimum size=4.6mm,
                    font=\sffamily\bfseries\fontsize{7}{8}\selectfont},
  rowTitle/.style={font=\sffamily\bfseries\fontsize{6.2}{7.4}\selectfont,
                   text=cConv, anchor=west},
  rowText/.style={font=\sffamily\fontsize{5.2}{6.2}\selectfont,
                  text=cGray!90!black, anchor=west},
  mathNode/.style={font=\fontsize{6}{7.2}\selectfont, text=cGray!90!black},
  arrowFlow/.style={-{Stealth[length=3pt,width=2.4pt]}, cGray!70!black,
                    line width=0.4pt},
]

% =========================================================================
% LEFT COLUMN — 3 numbered rows
% =========================================================================

% --- Row 1 : Experiment (log-log discharge plot) ---
\begin{scope}
  \path[rowFrame] (0.0, 7.4) rectangle (7.6, 10.6);
  \node[numCircle] at (0.55, 10.05) {1};
  \node[rowTitle] at (0.95, 10.20) {Experiment};
  \node[rowText]  at (0.95, 9.75)
    {Discharge power-law vs.\ time (log-log)};

  % Inset log-log plot via \LogLogPlot
  \LogLogPlot{3.4,7.85,6.9,10.25}
    {3.4,4.275,5.15,6.025,6.9}
    {3.4/{10^{0}}, 5.15/{10^{2}}, 6.9/{10^{4}}}
    {7.85,8.45,9.05,9.65,10.25}
    {7.85/{10^{-3}}, 9.05/{10^{-1}}, 10.25/{10^{1}}}
    {}{}
  % Axis labels removed (iter 2): row title already says "log-log",
  % tick labels carry magnitudes — own labels collided with ticks.

  % Power-law line (slope = -n) — solid blue
  \draw[cBlue, line width=0.85pt]
    (3.55, 10.10) -- (6.75, 8.05);
  \node[font=\sffamily\fontsize{5}{6}\selectfont, text=cBlue!85!black,
        anchor=west]
    at (5.40, 9.55) {Power-law};
  \node[font=\fontsize{5.4}{6.4}\selectfont, text=cBlue!85!black,
        anchor=west]
    at (5.40, 9.25) {$I = I_0\, t^{-n}$};
  \node[font=\sffamily\fontsize{4.8}{5.8}\selectfont, text=cBlue!85!black,
        anchor=west]
    at (5.40, 9.00) {(slope $=-n$)};

  % Debye reference — dashed gray exponential-ish curve (drawn as smooth poly-line)
  \draw[cGray!85, line width=0.6pt, dash pattern=on 1.6pt off 1.4pt]
    plot[smooth, samples=24, domain=0:1]
    ({3.55+\x*3.20}, {10.05 - 2.30*(1-exp(-3.0*\x))});
  \node[font=\sffamily\fontsize{5}{6}\selectfont, text=cGray!85,
        anchor=west]
    at (3.60, 8.30) {Debye reference};
  \node[font=\fontsize{5.4}{6.4}\selectfont, text=cGray!85,
        anchor=west]
    at (3.60, 8.05) {$I = I_0\, e^{-t/\tau}$};
\end{scope}

% --- Row 2 : Mathematical interpretation ---
\begin{scope}
  \path[rowFrame] (0.0, 3.7) rectangle (7.6, 6.9);
  \node[numCircle] at (0.55, 6.35) {2};
  \node[rowTitle] at (0.95, 6.50) {Mathematical};
  \node[rowTitle] at (0.95, 6.18) {interpretation};
  % Wide sub-text + chain sub-labels removed (iter 2):
  % they collided with the governing equation and with each other.
  % Row title + symbol-arrow chain + equation carry the message.

  % n -> tau_d -> g(E_t) chain
  \node[mathNode] at (1.30, 4.85) {$n$};
  \draw[arrowFlow] (1.65, 4.85) -- (2.75, 4.85);
  \node[mathNode] at (3.10, 4.85) {$\tau_d$};
  \draw[arrowFlow] (3.45, 4.85) -- (4.55, 4.85);
  \node[mathNode] at (4.95, 4.85) {$g(E_t)$};

  % Governing equation
  \node[font=\fontsize{6.5}{7.5}\selectfont, text=cGray!90!black]
    at (3.10, 5.55)
    {$\tau_d = \tau_0\, \exp\!\left(\dfrac{E_t}{k_B T}\right)$};

  % Inline two-peak g(E_t) sketch removed (iter 1, WARN-driven):
  % collisions with sub-labels and right-margin curves duplicate intent.
\end{scope}

% --- Row 3 : Molecular origin ---
\begin{scope}
  \path[rowFrame] (0.0, 0.0) rectangle (7.6, 3.2);
  \node[numCircle] at (0.55, 2.65) {3};
  \node[rowTitle] at (0.95, 2.80) {Molecular origin};
  \node[rowText]  at (0.95, 2.40)
    {Sulfur-rich polymer: chemical and physical origin of traps};

  % Polymer backbone — \WavyChain zigzag with S sites and trap markers.
  \WavyChain{1.10, 6.80, 1.55, zigzag,
             {1.85, 2.95, 4.05, 5.15, 6.25},
             {2.45, 4.55}}

  % Two captioned origin sub-bullets (chemical / physical)
  \path[draw=cAmber!75, line width=0.45pt, fill=cAmber!10,
        rounded corners=1.2mm]
    (0.85, 0.30) rectangle (3.85, 1.20);
  \node[font=\sffamily\bfseries\fontsize{5.2}{6.2}\selectfont,
        text=cAmber!85!black, anchor=west]
    at (0.95, 1.00) {chemical origin};
  \node[font=\sffamily\fontsize{4.6}{5.6}\selectfont,
        text=cGray!90!black, anchor=west, text width=2.85cm]
    at (0.95, 0.60) {electron-rich sulfur sites (polarizability, lone pairs)};

  \path[draw=cBlue!45!cRed!70, line width=0.45pt,
        fill=cBlue!45!cRed!10, rounded corners=1.2mm]
    (4.00, 0.30) rectangle (7.30, 1.20);
  \node[font=\sffamily\bfseries\fontsize{5.2}{6.2}\selectfont,
        text=cBlue!45!cRed!85!black, anchor=west]
    at (4.10, 1.00) {physical origin};
  \node[font=\sffamily\fontsize{4.6}{5.6}\selectfont,
        text=cGray!90!black, anchor=west, text width=3.10cm]
    at (4.10, 0.60) {conformational / free-volume heterogeneity};
\end{scope}

% =========================================================================
% CENTER — "Convergence" connector
% Hand-written: no published macro for vertical labeled connector boxes;
% briefing §5 calls for green-outlined accent that is not in BandDiagram or
% BellCurve scope.
% =========================================================================
\begin{scope}
  \path[draw=cConv, line width=0.9pt, fill=cConv!8,
        rounded corners=2.0mm]
    (7.95, 1.40) rectangle (8.95, 9.20);

  \node[rotate=90, font=\sffamily\bfseries\fontsize{6.5}{7.8}\selectfont,
        text=cConv]
    at (8.45, 7.80) {Convergence};
  \node[rotate=90, font=\sffamily\fontsize{5.8}{7.0}\selectfont,
        text=cConv]
    at (8.45, 3.20) {to a unified trap-depth picture};

  % Arrow head pointing right into the right-panel band diagram.
  \draw[-{Stealth[length=4.5pt, width=3.5pt]}, cConv, line width=0.9pt]
    (8.95, 5.30) -- (9.40, 5.30);
\end{scope}

% =========================================================================
% RIGHT PANEL — Unified band diagram + g(E_t) side curves
% =========================================================================

% Panel header
\node[font=\sffamily\bfseries\fontsize{7}{8.4}\selectfont,
      text=cGray!90!black, anchor=west]
    at (9.55, 10.30) {Unified trap-depth picture in a sulfur-rich polymer};

% \BandDiagram bbox: x1=9.55, y1=0.40, x2=13.10, y2=9.85
% CB at y=8.20, VB at y=2.00, Et reference at x=11.30
% Shallow trap dashes near CB; deep trap dashes near VB.
\BandDiagram
    {9.55, 0.40, 13.10, 9.85,
     8.20, 2.00, 11.30,
     {7.30, 7.00, 6.70},
     {3.40, 3.05}}

% External labels for shallow / deep trap clusters (briefing §6 verbatim).
\node[font=\sffamily\fontsize{5.6}{6.8}\selectfont, text=cAmber!90!black,
      anchor=west]
    at (12.10, 7.00) {Shallow};
\node[font=\sffamily\fontsize{5.6}{6.8}\selectfont, text=cAmber!90!black,
      anchor=west]
    at (12.10, 6.65) {traps};

\node[font=\sffamily\fontsize{5.6}{6.8}\selectfont,
      text=cBlue!45!cRed!85!black, anchor=west]
    at (12.10, 3.30) {Deep};
\node[font=\sffamily\fontsize{5.6}{6.8}\selectfont,
      text=cBlue!45!cRed!85!black, anchor=west]
    at (12.10, 2.95) {traps};

% =========================================================================
% RIGHT MARGIN — paired g(E_t) profiles aligned to shallow / deep regions.
% Two BellCurve calls, side orientation, vertically stacked.
% =========================================================================

% g(E_t) header
\node[font=\fontsize{6.2}{7.4}\selectfont, text=cGray!90!black]
    at (14.30, 9.50) {$g(E_t)$};

% Shallow-trap g(E_t): amber lobe at upper-right, aligned to shallow band.
\BellCurve[draw=cAmber!80!black, fill=cAmber!22, line width=0.7pt]
    {13.50, 5.80, 14.95, 8.30, side}
\draw[cGray!70, line width=0.32pt]
    (13.50, 5.80) -- (13.50, 8.30);     % E_t axis (vertical)
\draw[cGray!70, line width=0.32pt,
      -{Stealth[length=2pt,width=1.4pt]}]
    (13.50, 5.80) -- (15.10, 5.80);     % g(E_t) axis (horizontal)
\node[font=\fontsize{4.8}{5.8}\selectfont, text=cGray!85, anchor=north]
    at (15.10, 5.78) {$g(E_t)$};
\node[rotate=90, font=\fontsize{4.8}{5.8}\selectfont, text=cGray!85,
      anchor=south]
    at (13.40, 7.05) {$E_t$};

% Deep-trap g(E_t): violet lobe at lower-right, aligned to deep band.
\BellCurve[draw=cBlue!45!cRed!80!black, fill=cBlue!45!cRed!22,
           line width=0.7pt]
    {13.50, 1.20, 14.95, 4.20, side}
\draw[cGray!70, line width=0.32pt]
    (13.50, 1.20) -- (13.50, 4.20);
\draw[cGray!70, line width=0.32pt,
      -{Stealth[length=2pt,width=1.4pt]}]
    (13.50, 1.20) -- (15.10, 1.20);
\node[font=\fontsize{4.8}{5.8}\selectfont, text=cGray!85, anchor=north]
    at (15.10, 1.18) {$g(E_t)$};
\node[rotate=90, font=\fontsize{4.8}{5.8}\selectfont, text=cGray!85,
      anchor=south]
    at (13.40, 2.70) {$E_t$};

\end{tikzpicture}
\end{document}
